\section{AHSME 1955}
        \begin{enumerate}
        	\item (\textit{AHSME 1955}) Which one of the following is not equivalent to $0.000000375$? 


 $\textbf{(A)}\ 3.75\times 10^{-7}\qquad\textbf{(B)}\ 3\frac{3}{4}\times 10^{-7}\qquad\textbf{(C)}\ 375\times 10^{-9}\\ \textbf{(D)}\ \frac{3}{8}\times 10^{-7}\qquad\textbf{(E)}\ \frac{3}{80000000}$


 
 \textbf{Answer: }D



\textbf{Solution}

First of all, $0.000000375 = 3.75 \times 10^{-7}$ in scientific notation. This eliminates $\textbf{(A)}\ 3.75\times 10^{-7}$, $\textbf{(B)}\ 3\frac{3}{4}\times 10^{-7}$, and $\textbf{(C)}\ 375\times 10^{-9}$ immediately. $\textbf{(E)}\ \frac{3}{80000000}$ is a bit harder, but it can be rewritten as $\frac{3}{8 \times 10^{7}} = \frac{3}{8} \times 10^{-7}$, which is the same as $\textbf{(D)}$. Thus, the answer is either $\textbf{(D)}$ or $\textbf{(E)}$.  The original contest had this error in it, although the desired answer was D.  It is likely that they intended option E to be $\frac{3}{8000000}$, but misprinted it.

 


	\item (\textit{AHSME 1955}) The smaller angle between the hands of a clock at $12:25$ p.m. is: 


 $\textbf{(A)}\ 132^\circ 30'\qquad\textbf{(B)}\ 137^\circ 30'\qquad\textbf{(C)}\ 150^\circ\qquad\textbf{(D)}\ 137^\circ 32'\qquad\textbf{(E)}\ 137^\circ$


 
 \textbf{Answer: }B



\textbf{Solution}

At $12:25$, the minute hand is at $\frac{25}{60}\cdot 360^\circ$, or $150^\circ$.The hour hand moves $5$ 'minutes' every hour, or $\frac{5}{60}\cdot 360 = 30^\circ$ every hour. At $30^\circ$ every hour, the hour hand moves $\frac{1}{2}$ minutes on the clock every minute. At $12:25$, the hour hand is at $\frac{25}{2}^\circ$. Therefore, the angle between the hands is $150^\circ - \frac{25}{2}^\circ$, $137.5^\circ$, or $137^{\circ} 30^{'} \implies$ $\boxed{\mathrm{(B) 137^\circ 30'}}$.

 



\textbf{Solution 2}

Using the formula $\frac{|(60h-11m)|}{2}$, and plugging in 12 for h and 25 for m, we get the angle between to be 222.5$^{\circ}$. Since the problem wants the smaller angle, we do $360-222.5 = 137.5 \Rightarrow 137^{\circ}30'\Rightarrow \fbox{B}$

 


	\item (\textit{AHSME 1955}) If each number in a set of ten numbers is increased by $20$, the arithmetic mean (average) of the ten numbers: 


 $\textbf{(A)}\ \text{remains the same}\qquad\textbf{(B)}\ \text{is increased by 20}\qquad\textbf{(C)}\ \text{is increased by 200}\\ \textbf{(D)}\ \text{is increased by 10}\qquad\textbf{(E)}\ \text{is increased by 2}$


 
 \textbf{Answer: }B



\textbf{Solution 1}

Let the sum of the 10 numbers be x. The mean is then $\frac{x}{10}$.  Then, since you're adding 20 to each number, the new sum of the numbers is x+200, since there are 10 numbers. Then, the new mean is $\frac{x+200}{10}$, which simplifies to $\frac{x}{10}+20$, which is $\fbox{B}$.

 


	\item (\textit{AHSME 1955}) The equality $\frac{1}{x-1}=\frac{2}{x-2}$ is satisfied by: 


 $\textbf{(A)}\ \text{no real values of }x\qquad\textbf{(B)}\ \text{either }x=1\text{ or }x=2\qquad\textbf{(C)}\ \text{only }x=1\\ \textbf{(D)}\ \text{only }x=2\qquad\textbf{(E)}\ \text{only }x=0$


 
 \textbf{Answer: }E



\textbf{Solution 1}

From the equality, $\frac{1}{x-1}=\frac{2}{x-2}$, we get  ${(x-1)}\times2={(x-2)}\times1$.

 Solving this, we get, ${2x-2}={x-2}$.

 Thus, the answer is $\fbox{{\bf(E)} \text{only} x = 0}$.

 Solution by awesomechoco

 



\textbf{Solution 2 (answer choices)}

Looking at the answer choices, we can see B, C, and D are incorrect because $x = 1$ and $x = 2$ result in division by $0$.

 Plugging in $0$ gives:

 $\frac{1}{0 - 1} = \frac{1}{-1} = -1$

 $\frac{2}{0 - 2} = \frac{2}{-2} = -1$

 These two are equal, so the answer is $\fbox{{\bf(E)} \text{only} x = 0}$.

 


	\item (\textit{AHSME 1955}) $5y$ varies inversely as the square of $x$. When $y=16, x=1$. When $x=8, y$ equals: 


 $\textbf{(A)}\ 2 \qquad \textbf{(B)}\ 128 \qquad \textbf{(C)}\ 64 \qquad \textbf{(D)}\ \frac{1}{4} \qquad \textbf{(E)}\ 1024$


 
 \textbf{Answer: }D



\textbf{Solution}

An inverse variation can be expressed in the form $xy = n$, where $n$ is any number (except perhaps zero). Since $5y$ varies inversely with the square of $x$, this particular one will be $5yx^2 = n$. 

 We can plug in $16$ for $y$ and $1$ for $x$, which makes $n=80$. The equation is now $5yx^2 = 80$.

 When $x = 8$, we can solve for $y$ using the equation $320y = 80$, which makes $y$ $\boxed{\textbf{(D)} \frac{1}{4}}$.

 


	\item (\textit{AHSME 1955}) A merchant buys a number of oranges at $3$ for $10$ cents and an equal number at $5$ for $20$ cents. To "break even" he must sell all at: 


 $\textbf{(A)}\ \text{8 for 30 cents}\qquad\textbf{(B)}\ \text{3 for 11 cents}\qquad\textbf{(C)}\ \text{5 for 18 cents}\\ \textbf{(D)}\ \text{11 for 40 cents}\qquad\textbf{(E)}\ \text{13 for 50 cents}$


 
 \textbf{Answer: }B



\textbf{Solution}

Since we are buying at $3$ for $10$ cents and $5$ for $20$ cents, let's assume that together, we are buying 15 oranges.That means that we are getting a total of $30$ oranges for $(10\times5) + (20\times3)$ cents.That comes to a total of $30$ oranges for $110$ cents. $110/30$ = $11/3$. This leads us to $3$ for $11$ cents which is $\boxed{B}$ and we are done.

 -Brudder

 Edited by Andrew Lu  (originally answer was C)

 


	\item (\textit{AHSME 1955}) If a worker receives a $20$\% cut in wages, he may regain his original pay exactly by obtaining a raise of:


 $\textbf{(A)}\ \text{20\%}\qquad\textbf{(B)}\ \text{25\%}\qquad\textbf{(C)}\ 22\frac{1}{2}\text{\%}\qquad\textbf{(D)}\ \textdollar{20}\qquad\textbf{(E)}\ \textdollar{25}$


 
 \textbf{Answer: }B



\textbf{Solution}

Since the worker receives a $20$\% cut in wages, his present wage would be $\frac{4}{5}$ of the original wage before the cut. To regain his original pay he would have to obtain a raise of $1$ $\div$ $\frac{4}{5}$ = $1$$\frac{1}{4}$.

 Therefore, the worker would have to get a $\fbox{{\bf(B)} 25\%}$ raise to acquire his original pay.

 Solution by awesomechoco

 


	\item (\textit{AHSME 1955}) The graph of $x^2-4y^2=0$: 


 $\textbf{(A)}\ \text{is a hyperbola intersecting only the }x\text{-axis}\\ \textbf{(B)}\ \text{is a hyperbola intersecting only the }y\text{-axis}\\ \textbf{(C)}\ \text{is a hyperbola intersecting neither axis}\\ \textbf{(D)}\ \text{is a pair of straight lines}\\ \textbf{(E)}\ \text{does not exist}$


 
 \textbf{Answer: }D



\textbf{Solution}

By difference of squares, we can rewrite the equation as $(x-2y)(x+2y) = 0$, which is just the union of the two lines $x - 2y = 0$ and $x + 2y = 0$. Therefore, our answer is $\boxed{\textbf{(D)}}$, and we are done.

 


	\item (\textit{AHSME 1955}) A circle is inscribed in a triangle with sides $8, 15$, and $17$. The radius of the circle is: 


 $\textbf{(A)}\ 6 \qquad \textbf{(B)}\ 2 \qquad \textbf{(C)}\ 5 \qquad \textbf{(D)}\ 3 \qquad \textbf{(E)}\ 7$


 
 \textbf{Answer: }D



\textbf{Solution}

We know that $A = sr$, where $A$ is the triangle's area, $s$ its semiperimeter, and $r$ its inradius. Since this particular triangle is a right triangle (which we can verify by the Pythagorean theorem), the area is half of $8*15 = 120$, and the semiperimeter is half of $8 + 15 + 17 = 40$. Therefore, the inradius is $\frac{120}{40} = 3$, so our answer is $\boxed{\textbf{(D)}}$ and we are done.

 


 


	\item (\textit{AHSME 1955}) How many hours does it take a train traveling at an average rate of 40 mph between stops to travel a miles it makes n stops of m minutes each? 


 $\textbf{(A)}\ \frac{3a+2mn}{120}\qquad\textbf{(B)}\ 3a+2mn\qquad\textbf{(C)}\ \frac{3a+2mn}{12}\qquad\textbf{(D)}\ \frac{a+mn}{40}\qquad\textbf{(E)}\ \frac{a+40mn}{40}$


 
 \textbf{Answer: }A



\textbf{Solution}

The train will take $\frac{a}{40}$ hours to travel $a$ miles, and it takes $\frac{nm}{60}$. The LCM of $40$ and $60$ is $120$, which allows for the addition of the fractions $\frac{3a}{120}$ and $\frac{2mn}{120}$. The end result is $\boxed{\textbf{(A)}\frac{3a+2mn}{120}}$.

 


	\item (\textit{AHSME 1955}) The negation of the statement "No slow learners attend this school" is: 


 $\textbf{(A)}\ \text{All slow learners attend this school}\\ \textbf{(B)}\ \text{All slow learners do not attend this school}\\ \textbf{(C)}\ \text{Some slow learners attend this school}\\ \textbf{(D)}\ \text{Some slow learners do not attend this school}\\ \textbf{(E)}\ \text{No slow learners do not attend this school}$


 
 \textbf{Answer: }C



\textbf{Solution}

The negation of the statement is basically the opposite of the statement, which is $\boxed{\textbf{(C) some slow learners attend this school.}}$

 


	\item (\textit{AHSME 1955}) The solution of $\sqrt{5x-1}+\sqrt{x-1}=2$ is: 


 $\textbf{(A)}\ x=2,x=1\qquad\textbf{(B)}\ x=\frac{2}{3}\qquad\textbf{(C)}\ x=2\qquad\textbf{(D)}\ x=1\qquad\textbf{(E)}\ x=0$


 
 \textbf{Answer: }D



\textbf{Solution 1}

First, square both sides.  This gives us

 
 \[\sqrt{5x-1}^2+2\cdot\sqrt{5x-1}\cdot\sqrt{x-1}+\sqrt{x-1}^2=4 \Longrightarrow 5x-1+2\cdot\sqrt{(5x-1)\cdot(x-1)}+x-1=4 \Longrightarrow 2\cdot\sqrt{5x^2-6x+1}+6x-2=4\]Then, adding $-6x$ to both sides gives us

 
 \[2\cdot\sqrt{5x^2-6x+1}+6x-2-6x=4-6x \Longrightarrow 2\cdot\sqrt{5x^2-6x+1}-2 =-6x+4\]After that, adding $2$ to both sides will give us

 
 \[2\cdot\sqrt{5x^2-6x+1}-2+2=-6x+4+2 \Longrightarrow 2\cdot\sqrt{5x^2-6x+1}=-6x+6\]Next, we divide both sides by 2 which gives us

 
 \[\frac{2\cdot\sqrt{5x^2-6x+1}-2+2}{2}=\frac{-6x+4+2}{2} \Longrightarrow \sqrt{5x^2-6x+1}=-3x+3\]Finally, solving the equation, we get 

 
 \[5x^2-6x+1=(-3x+3)^2 \Longrightarrow 5x^2-6x+1=9x^2-18x+9\] 
 \[\Longrightarrow 5x^2-6x+1-(9x^2-18x+9)=9x^2-18x+9-(9x^2-18x+9)\]
 \[\Longrightarrow -4x^2+12x-8=0 \Longrightarrow -4(x-1)(x-2)=0\]
 \[\Longrightarrow x-1=0 \text{or}\ x-2=0 \Longrightarrow x=1 \text{or}\ x=2\]Plugging 1 and 2 into the original equation, $\sqrt{5x-1}+\sqrt{x-1}=2$, we see that when $x=1$

 
 \[\sqrt{5x-1}+\sqrt{x-1}=2 \Longrightarrow \sqrt{5\cdot1-1}+\sqrt{1-1}=2 \Longrightarrow \sqrt4+\sqrt0=2 \Longrightarrow 2+0=2 \Longrightarrow 2=2\]the equation is true. On the other hand, we note that when $x=2$

 
 \[\sqrt{5x-1}+\sqrt{x-1}=2 \Longrightarrow \sqrt{5\cdot2-1}+\sqrt{2-1}=2 \Longrightarrow \sqrt9+\sqrt1=2 \Longrightarrow 3+1=2 \Longrightarrow 4=0\]the equation is false. Therefore the answer is $\boxed{{\textbf{(D) }} x=1}$.

 



\textbf{Solution 2}

Let us test the answer choices, for it is in this case simpler and quicker. $x=0$ and $x=2/3$ obviously doesn't work, since square roots of negative numbers are not real. $x=2$ doesn't work either, because the first term is already bigger than 2, and the second term cannot be negative. By the process of elimination, the answer is $\boxed{{\textbf{(D) }} x=1}$.

 


	\item (\textit{AHSME 1955}) The fraction $\frac{a^{-4}-b^{-4}}{a^{-2}-b^{-2}}$ is equal to: 


 $\textbf{(A)}\ a^{-6}-b^{-6}\qquad\textbf{(B)}\ a^{-2}-b^{-2}\qquad\textbf{(C)}\ a^{-2}+b^{-2}\\ \textbf{(D)}\ a^2+b^2\qquad\textbf{(E)}\ a^2-b^2$


 
 \textbf{Answer: }C



\textbf{Solution}

By the difference of squares property, $a^{-4} - b^{-4}$ is equivalent to $(a^{-2} + b^{-2})(a^{-2} - b^{-2})$. This means the fraction is now equal to $\frac{(a^{-2} + b^{-2})(a^{-2} - b^{-2})}{a^{-2} - b^{-2}}$, which simplifies to $\textbf{(C)}\ a^{-2}+b^{-2}$ 

 


	\item (\textit{AHSME 1955}) The length of rectangle $R$ is $10$\% more than the side of square $S$. The width of the rectangle is $10$\% less than the side of the square. 
The ratio of the areas, $R:S$, is: 


 $\textbf{(A)}\ 99: 100\qquad\textbf{(B)}\ 101: 100\qquad\textbf{(C)}\ 1: 1\qquad\textbf{(D)}\ 199: 200\qquad\textbf{(E)}\ 201: 200$


 
 \textbf{Answer: }A



\textbf{Solution}

Let each of the square's sides be $x$. The dimensions of the rectangle can be expressed as $1.1x$ and $0.9x$. Therefore, the area of the rectangle is $0.99x^2$, while the square has an area of $x^2$. The ratio of $R : S$ can be defined as $0.99x^2 : x^2$, which ultimately leads to $\textbf{(A) } 99 : 100$

 


	\item (\textit{AHSME 1955}) The ratio of the areas of two concentric circles is $1: 3$. If the radius of the smaller is $r$, then the difference between the 
radii is best approximated by: 


 $\textbf{(A)}\ 0.41r \qquad \textbf{(B)}\ 0.73 \qquad \textbf{(C)}\ 0.75 \qquad \textbf{(D)}\ 0.73r \qquad \textbf{(E)}\ 0.75r$


 
 \textbf{Answer: }D



\textbf{Solution}

Observe that there are 3 feet in 1 yard, 9 square feet in 1 square yard, and 27 cubic feet in 1 cubic yard.

 This means that if the ratio is $x$ in 1D, it is $x^2$ in 2D, and it is $x^3$ in 3D.

 We can apply this thinking into our current question, which would have $1:3$ in the 2D-plane. Since the radius is 1D, we would have to square root both sides to get $1:\sqrt{3}$

 If we plug in $r$ for the smaller radius, we get $r:\sqrt{3}r$ The difference between the two sides would be (approximately) $1.732r - r = \textbf{(D)} 0.73r$

 


	\item (\textit{AHSME 1955}) The value of $\frac{3}{a+b}$ when $a=4$ and $b=-4$ is: 


 $\textbf{(A)}\ 3\qquad\textbf{(B)}\ \frac{3}{8}\qquad\textbf{(C)}\ 0\qquad\textbf{(D)}\ \text{any finite number}\qquad\textbf{(E)}\ \text{meaningless}$


 
 \textbf{Answer: }E



\textbf{Solution}

As the denominator is $4 - (-4) = 0$, the value of the fraction $\frac{3}{4 - (-4)}$ is $\textbf{(E)}\ \text{meaningless}$

 


	\item (\textit{AHSME 1955}) If $\log x-5 \log 3=-2$, then $x$ equals: 


 $\textbf{(A)}\ 1.25\qquad\textbf{(B)}\ 0.81\qquad\textbf{(C)}\ 2.43\qquad\textbf{(D)}\ 0.8\qquad\textbf{(E)}\ \text{either 0.8 or 1.25}$


 
 \textbf{Answer: }C



\textbf{Solution}

\textbf{Definitions and Properties} $\log n$ is defined as the value of $x$ that satisfies the equation $n = 10^x$. Note that other bases can be applied as well, so $\log_b n$ would be defined as the answer to $n = b^x$

 $y \log n = \log n^y$

 $\log x - \log y = \log (x / y)$

 Learn more properties here: [\href{https://www.rapidtables.com/math/algebra/logarithm/Logarithm\_Rules.html}{[1]}]

 \textbf{Solving the Equation} $\log x-5 \log 3=-2$

 $\log x- \log 3^5=-2$

 $\log x- \log 243 =-2$

 $\log x / 243 = -2$

 $x/243 = 10^{-2}$

 $x=\frac{243}{100}$$x=\boxed{\textbf{(C) }2.43}$

 


	\item (\textit{AHSME 1955}) The discriminant of the equation $x^2+2x\sqrt{3}+3=0$ is zero. Hence, its roots are: 


 $\textbf{(A)}\ \text{real and equal}\qquad\textbf{(B)}\ \text{rational and equal}\qquad\textbf{(C)}\ \text{rational and unequal}\\ \textbf{(D)}\ \text{irrational and unequal}\qquad\textbf{(E)}\ \text{imaginary}$


 
 \textbf{Answer: }A



\textbf{Solution}

Since the discriminant is zero, there is one distinct root, or, relevant to this question, two equal roots. The fact that only one solution exists means that the equation can be simplified into $(x + \sqrt{3})^2=0$. The distinct solution\textasciicircum{} to the equation, as we can clearly see, is $-\sqrt{3}$, which is $\textbf{(A)} \text{real and equal}$

 \textasciicircum{}The problem pretends to have two solutions.

 


	\item (\textit{AHSME 1955}) Two numbers whose sum is $6$ and the absolute value of whose difference is $8$ are roots of the equation: 


 $\textbf{(A)}\ x^2-6x+7=0\qquad\textbf{(B)}\ x^2-6x-7=0\qquad\textbf{(C)}\ x^2+6x-8=0\\ \textbf{(D)}\ x^2-6x+8=0\qquad\textbf{(E)}\ x^2+6x-7=0$


 
 \textbf{Answer: }B



\textbf{Solution}

The first two hints can be expressed as the following system of equations:
 \[\begin{cases} (1) & a + b = 6 \\ (2) & a - b = 8 \end{cases}\]From this, we can clearly see that $a = 7$, and that $b = -1$.

 Since quadratic equations can generally be expressed in the form of $(x - a)(x - b) = 0$, where a and b are roots, the correct quadratic, once factored, would look like $(x - 7)(x + 1) = 0$

 Expanding the above equation gets us $\textbf{(B)} x^2 - 6x - 7 = 0$

 



\textbf{Solution 2}

Let the roots of the equation be $x$ and $y$. Therefore, we can set up a system of equation:
 \[x+y=6\]
 \[|x-y|=8\]Therefore, we get $x=7$ and $y=-1$. So, $(x-7)(x+1)=\boxed{x^2-6x-7}$

 - kante314 -

 


	\item (\textit{AHSME 1955}) The expression $\sqrt{25-t^2}+5$ equals zero for: 


 $\textbf{(A)}\ \text{no real or imaginary values of }t\qquad\textbf{(B)}\ \text{no real values of }t\text{ only}\\ \textbf{(C)}\ \text{no imaginary values of }t\text{ only}\qquad\textbf{(D)}\ t=0\qquad\textbf{(E)}\ t=\pm 5$


 
 \textbf{Answer: }A



\textbf{Solution}

We can make the following equation:  $\sqrt{25-t^2}+5 = 0$, which means that $\sqrt{25-t^2} = -5$ Square both sides, and we get $25 - t^2 = 25$, which would mean that $t = 0$, but that is an extraneous solution and doesn't work in the original equation.

 Therefore, the correct answer is $\textbf{(A)}$

 Note: When talking about square roots, it generally means the positive square root, or the "principal root."

 


	\item (\textit{AHSME 1955}) Represent the hypotenuse of a right triangle by $c$ and the area by $A$. The altitude on the hypotenuse is: 


 $\textbf{(A)}\ \frac{A}{c}\qquad\textbf{(B)}\ \frac{2A}{c}\qquad\textbf{(C)}\ \frac{A}{2c}\qquad\textbf{(D)}\ \frac{A^2}{c}\qquad\textbf{(E)}\ \frac{A}{c^2}$


 
 \textbf{Answer: }B



\textbf{Solution}


                    \begin{figure}[H]
                        centering
                        \includegraphics[width=0.5\textwidth]{Solution/AHSME_I_1955_21_Solution_Figure_1.png}
                    \end{figure}
                    Given that the area of the triangle is $A$, and the formula for the area of a triangle is $\frac{bh}{2}=A$, we can replace $b$ (the base) and $h$ (the height) with $c$ (the hypotenuse) and $a$ (the altitude), we can rearrrange as follows:
 \[\frac{c*a}{2}=A\]
 \[c*a=2A\]
 \[a=\textbf{(B)} \frac{2A}{c}\]

 


	\item (\textit{AHSME 1955}) On a $\textdollar{10000}$ order a merchant has a choice between three successive discounts of $20$\%, $20$\%, and $10$\% and 
three successive discounts of $40$\%, $5$\%, and $5$\%. By choosing the better offer, he can save: 


 $\textbf{(A)}\ \text{nothing at all}\qquad\textbf{(B)}\ \$440\qquad\textbf{(C)}\ \$330\qquad\textbf{(D)}\ \$345\qquad\textbf{(E)}\ \$360$


 
 \textbf{Answer: }D



\textbf{Solution}

In order to solve this problem, we can simply try both paths and finding the positive difference between the two.

 The first path goes $10,000 * 0.8 = 8,000 * 0.8 = 6400 * 0.9 = \textbf{5760}$

 The second path goes $10,000 * 0.6 = 6,000 * 0.95 \text{ or } (1 - 0.05) = 5700 * 0.95 = \textbf{5415}$

 The positive difference between the offers is therefore $5760 - 5415 = \fbox{\textbf{(D)} \textdollar{345}}$

 


	\item (\textit{AHSME 1955}) In checking the petty cash a clerk counts $q$ quarters, $d$ dimes, $n$ nickels, and $c$ cents. 
Later he discovers that $x$ of the nickels were counted as quarters and $x$ of the dimes were counted as cents. 
To correct the total obtained the clerk must: 


 $\textbf{(A)}\ \text{make no correction}\qquad\textbf{(B)}\ \text{subtract 11 cents}\qquad\textbf{(C)}\ \text{subtract 11}x\text{ cents}\\ \textbf{(D)}\ \text{add 11}x\text{ cents}\qquad\textbf{(E)}\ \text{add }x\text{ cents}$


 
 \textbf{Answer: }C



\textbf{Solution}

If the clerk mistook $x$ nickels as quarters, then every mistake inflates the total by $20$ cents. In order to correct this, we have to subtract $20$ cents $x$ times, for a total of $20x$ cents. We can do the same for the $x$ dimes that were turned into pennies (or cents). This exchange would increase the total value of the coins by $9x$.

 As a total, you get $-20x + 9x = -11x$. In order to correct the amount, the clerk should $\boxed{\text{subtract } 11x \text{ cents}}$.

 


	\item (\textit{AHSME 1955}) The function $4x^2-12x-1$: 


 $\textbf{(A)}\ \text{always increases as }x\text{ increases}\\ \textbf{(B)}\ \text{always decreases as }x\text{ decreases to 1}\\ \textbf{(C)}\ \text{cannot equal 0}\\ \textbf{(D)}\ \text{has a maximum value when }x\text{ is negative}\\ \textbf{(E)}\ \text{has a minimum value of-10}$


 
 \textbf{Answer: }E



\textbf{Solution}

We can use the process of elimination to narrow down the field substantially:

 $\textbf{(A)}\ \text{always increases as } x\text{ increases}$ is wrong due to the quadratic nature of the function.

 $\textbf{(B)}\ \text{always decreases as } x\text{ decreases to 1}$ is wrong due to the vertex being on the line $x = 1.5$.  $f(x)$ would decrease all the way to $x = 1.5$, but increase from there.

 $\textbf{(C)}\ \text{cannot equal 0}$ is wrong due to the \href{https://artofproblemsolving.com/wiki/index.php?title=Discriminant}{discriminant} being greater than zero.

 $\textbf{(D)}\ \text{has a maximum value when}\ x \text{ is negative}$ is wrong due to the function not having a maximum value (the function opens up, so the vertex is the minimum value).

 Therefore, $\textbf{(E)}$ is the correct answer (this can be verified by plugging in 1.5 (the x-coordinate of the vertex) in).

 


	\item (\textit{AHSME 1955}) One of the factors of $x^4+2x^2+9$ is: 


 $\textbf{(A)}\ x^2+3\qquad\textbf{(B)}\ x+1\qquad\textbf{(C)}\ x^2-3\qquad\textbf{(D)}\ x^2-2x-3\qquad\textbf{(E)}\ \text{none of these}$


 
 \textbf{Answer: }E



\textbf{Solution}

We can test each of the answer choices by using polynomial division. 

 $x^2 + 3$ leaves behind a remainder, and so does $x^2 - 3$.

 In addition, $x + 1$ also fails the test, and that takes down $x^2 - 2x - 3$, which can be expressed as $(x + 1)(x - 3)$. That leaves $\boxed{(\textbf{E})}$

 



\textbf{Solution 2 (direct factorization)}

Notice the leading and constant terms are begging us to create a binomial. So
 \[x^4 + 2x^2 + 9 = (x^4 + 6x^2 + 9) - 4x^2 = (x^2 + 3)^2 - (2x)^2 = (x^2 + 2x + 3)(x^2 - 2x + 3),\]where both quadratics are irreducible (over the field of real numbers).Hence none of the given options is a factor. So the answer is $\boxed{(\textbf{E})}$

 


	\item (\textit{AHSME 1955}) Mr. A owns a house worth $\textdollar{10000}$. He sells it to Mr. $B$ at $10$\% profit. Mr. $B$ sells the house back to Mr. $A$ at a $10$\% loss. Then: 


 $\textbf{(A)}\ \text{Mr. A comes out even}\qquad\textbf{(B)}\ \text{Mr. A makes } \textdollar{ 100}\qquad\textbf{(C)}\ \text{Mr. A makes } \textdollar{ 1000}\\ \textbf{(D)}\ \text{Mr. B loses } \textdollar{ 100}\qquad\textbf{(E)}\ \text{none of the above is correct}$


 
 \textbf{Answer: }E



\textbf{Solution}

Assume that Mr. $A$ has $\textdollar{10000}$ to start with (the house) and Mr. $B$ has $\textdollar{11000}$ 

 $A$ sells the house to $B$ for $\textdollar{11000}$, leaving Mr. $A$ with $\textdollar{11000}$ and Mr. $B$ with $\textdollar{10000}$. 

 Then, Mr. $B$ sells the house back to Mr. $A$ for $\textdollar{9000}$, which leaves $A$ with $\textdollar{12000}$ and $B$ with $\textdollar{9000}$.

 We can see that $A$ gained $\textdollar{2000}$ out of Mr. $B$'s expense, but none of the options says this. We can therefore conclude that the correct answer is $\boxed{\textbf{(C)}}$.

 


	\item (\textit{AHSME 1955}) If $r$ and $s$ are the roots of $x^2-px+q=0$, then $r^2+s^2$ equals: 


 $\textbf{(A)}\ p^2+2q\qquad\textbf{(B)}\ p^2-2q\qquad\textbf{(C)}\ p^2+q^2\qquad\textbf{(D)}\ p^2-q^2\qquad\textbf{(E)}\ p^2$


 
 \textbf{Answer: }B



\textbf{Solution}

We can write $r^2+s^2$ in terms of the sum of the roots $(r+s)$ and the products of the roots $(rs):$
 \[r^2 + s^2 = (r+s)^2 - 2rs = p^2 - 2q\]The answer is $\boxed{\textbf{(B)}}.$

 


	\item (\textit{AHSME 1955}) On the same set of axes are drawn the graph of $y=ax^2+bx+c$ and the graph of the equation obtained by replacing $x$ by $-x$ in the given equation. 
If $b \neq 0$ and $c \neq 0$ these two graphs intersect: 


 $\textbf{(A)}\ \text{in two points, one on the x-axis and one on the y-axis}\\ \textbf{(B)}\ \text{in one point located on neither axis}\\ \textbf{(C)}\ \text{only at the origin}\\ \textbf{(D)}\ \text{in one point on the x-axis}\\ \textbf{(E)}\ \text{in one point on the y-axis}$


 
 \textbf{Answer: }E



\textbf{Solution}

Replacing $x$ with $-x,$ we have the graph of $y = a(-x)^2 + b(-x) + c = ax^2 - bx+c.$

 We can plug in simple values of $a, b,$ and $c$ for convenient drawing. For example, we can graph $x^2+2x+1$ and $x^2-2x+1.$ We see that the parabolas intersect at $(0, 1),$ which is on the y-axis.

 The answer is $\boxed{\textbf{(E)}}.$

 


	\item (\textit{AHSME 1955}) In the figure, $PA$ is tangent to semicircle $SAR$; $PB$ is tangent to semicircle $RBT$; $SRT$ is a straight line; 
the arcs are indicated in the figure. $\angle APB$ is measured by: 


 
                    \begin{figure}[H]
                        centering
                        \includegraphics[width=0.5\textwidth]{Problem/AHSME_I_1955_29_Problem_Figure_1.png}
                    \end{figure}
                    


 $\textbf{(A)}\ \frac{1}{2}(a-b)\qquad\textbf{(B)}\ \frac{1}{2}(a+b)\qquad\textbf{(C)}\ (c-a)-(d-b)\qquad\textbf{(D)}\ a-b\qquad\textbf{(E)}\ a+b$


 
 \textbf{Answer: }E



\textbf{Solution}


                    \begin{figure}[H]
                        centering
                        \includegraphics[width=0.5\textwidth]{Solution/AHSME_I_1955_29_Solution_Figure_1.png}
                    \end{figure}
                    

 


	\item (\textit{AHSME 1955}) Each of the equations $3x^2-2=25, (2x-1)^2=(x-1)^2, \sqrt{x^2-7}=\sqrt{x-1}$ has: 


 $\textbf{(A)}\ \text{two integral roots}\qquad\textbf{(B)}\ \text{no root greater than 3}\qquad\textbf{(C)}\ \text{no root zero}\\ \textbf{(D)}\ \text{only one root}\qquad\textbf{(E)}\ \text{one negative root and one positive root}$


 
 \textbf{Answer: }B



\textbf{Solution}

Since the question asks us about the unifying characteristic of all three equations' roots, we have to first determine them.

 $3x^2-2 = 25$ can be rewritten as $3x^2 - 27 = 0$, which gives the following roots $+3$ and $-3$.

 $(2x-1)^2 = (x-1)^2$ can be expanded to $4x^2-4x+1=x^2-2x+1$, which in turn leads to $3x^2-2x=0$. The roots here are $0$ and $\frac{2}{3}$.

 $\sqrt{x^2-7}=\sqrt{x-1}$, when squared, also turns into a quadratic equation: $x^2 - x - 6 = 0$. Binomial factoring gives us the roots $-2$ and $3$.

 We can clearly see that, between all of the equations, there is $\boxed{\textbf{(B)} \text{no root greater than 3}}$.

 Note: There are probably extraneous solutions somewhere, but that does not affect the solution.

 


	\item (\textit{AHSME 1955}) An equilateral triangle whose side is $2$ is divided into a triangle and a trapezoid by a line drawn parallel to one of its sides. 
If the area of the trapezoid equals one-half of the area of the original triangle, the length of the median of the trapezoid is: 


 $\textbf{(A)}\ \frac{\sqrt{6}}{2}\qquad\textbf{(B)}\ \sqrt{2}\qquad\textbf{(C)}\ 2+\sqrt{2}\qquad\textbf{(D)}\ \frac{2+\sqrt{2}}{2}\qquad\textbf{(E)}\ \frac{2\sqrt{3}-\sqrt{6}}{2}$


 
 \textbf{Answer: }D



\textbf{Solution}

The area of the large equilateral triangle is $2^2*\frac{\sqrt{3}}{4}=\sqrt{3}$ square units, so the smaller triangle is $\frac{\sqrt{3}}{2}$. The side length of the smaller length $s$ is $\sqrt{\frac{4}{\sqrt{3}}*\frac{\sqrt{3}}{2}} = \sqrt{2}$. The median of the trapezoid can be determined by the average of the two bases of the trapezoids, so it works out to be $\boxed{\textbf{(D)}\frac{2+\sqrt{2}}{2}}$ units long.

 


	\item (\textit{AHSME 1955}) If the discriminant of $ax^2+2bx+c=0$ is zero, then another true statement about $a, b$, and $c$ is that: 


 $\textbf{(A)}\ \text{they form an arithmetic progression}\\ \textbf{(B)}\ \text{they form a geometric progression}\\ \textbf{(C)}\ \text{they are unequal}\\ \textbf{(D)}\ \text{they are all negative numbers}\\ \textbf{(E)}\ \text{only b is negative and a and c are positive}$


 
 \textbf{Answer: }B



\textbf{Solution}

The discriminant of a quadratic is

 
 \[\Delta = b^2 - 4ac = (2b)^2 - 4ac = 4b^2 - 4ac = 0.\]We know that $b^2 = ac,$ or $\frac{b}{a} = \frac{c}{b}$ so $a, b, c$ form a geometric progression.

 
$\boxed{\text{(B)}}$

 


	\item (\textit{AHSME 1955}) Henry starts a trip when the hands of the clock are together between $8$ a.m. and $9$ a.m. 
He arrives at his destination between $2$ p.m. and $3$ p.m. when the hands of the clock are exactly $180^\circ$ apart. The trip takes: 


 $\textbf{(A)}\ \text{6 hr.}\qquad\textbf{(B)}\ \text{6 hr.   } 43\frac{7}{11} \text{min.}\qquad\textbf{(C)}\ \text{5 hr.   } 16\frac{4}{11} \text{min.}\qquad\textbf{(D)}\ \text{6 hr. 30 min.}\qquad\textbf{(E)}\ \text{none of these}$


 
 \textbf{Answer: }A


	\item (\textit{AHSME 1955}) A $6$-inch and $18$-inch diameter pole are placed together and bound together with wire. 
The length of the shortest wire that will go around them is: 


 $\textbf{(A)}\ 12\sqrt{3}+16\pi\qquad\textbf{(B)}\ 12\sqrt{3}+7\pi\qquad\textbf{(C)}\ 12\sqrt{3}+14\pi\\ \textbf{(D)}\ 12+15\pi\qquad\textbf{(E)}\ 24\pi$


 
 \textbf{Answer: }C



\textbf{Solution}


                    \begin{figure}[H]
                        centering
                        \includegraphics[width=0.5\textwidth]{Solution/AHSME_I_1955_34_Solution_Figure_1.png}
                    \end{figure}
                    
 \[\text{A diagram of the problem...aside from a few blips.}\]We can separate the total length of the wire into two sections: the straight and the curve.

 In order to find out the length of the straight, we can combine the two radii (halved because $6$ and $18$ are diameters) to be the hypotenuse and the difference between said radii to be the short leg. With this, we can establish a $30$-$60$-$90$ triangle, which means that the remaining leg must be $6\sqrt{3}$ This is the length of one part of the straight. Since there are two parts, the total length is $12\sqrt{3}$.

 Next, we find the curved part of the wire. We must first find the arc length of the curves, and the angle measurements to find the lengths. We can do this by doing two different things for two different circles. For the radius $3$ circle, we can subtract two right angles and two $30$ degree angles from the total $360$, which makes for a total of $360 - 240 = 120$ degree measure. That means that the curved section on the radius $3$ is $\frac{1}{3}6\pi=2\pi$. For the bigger circle, we can simply subtract twice the $60$ degree angle from the circle to get $360 - 120 = 240$ degrees. When plugged into the formula, then we get $\frac{2}{3}18\pi=12\pi$.

 Add everything up and you get $\boxed{\textbf{(C)} 12\sqrt{3} + 14\pi}$

 


	\item (\textit{AHSME 1955}) Three boys agree to divide a bag of marbles in the following manner. The first boy takes one more than half the marbles. 
The second takes a third of the number remaining. The third boy finds that he is left with twice as many marbles as the second boy. 
The original number of marbles: 


 $\textbf{(A)}\ \text{is none of the following}\qquad\textbf{(B)}\ \text{cannot be determined from the given data}\\ \textbf{(C)}\ \text{is 20 or 26}\qquad\textbf{(D)}\ \text{is 14 or 32}\qquad\textbf{(E)}\ \text{is 8 or 38}$


 
 \textbf{Answer: }B



\textbf{Solution}


                    \begin{figure}[H]
                        centering
                        
\includegraphics[width=0.5\textwidth]{Solution/AHSME_I_1955_35_Solution_Figure_1.png}
                    \end{figure}
                    The line that represents the marbles is split into three sections: $3x+1,  x, \text{and } 2x$. Since the amount of marbles each boy relies on the variable $x$, which cannot be solved for, the original amount of marbles $\boxed{\textbf{(B) } \text{cannot be determined from the given data}}$.

 


	\item (\textit{AHSME 1955}) A cylindrical oil tank, lying horizontally, has an interior length of $10$ feet and an interior diameter of $6$ feet. 
If the rectangular surface of the oil has an area of $40$ square feet, the depth of the oil is: 


 $\textbf{(A)}\ \sqrt{5}\qquad\textbf{(B)}\ 2\sqrt{5}\qquad\textbf{(C)}\ 3-\sqrt{5}\qquad\textbf{(D)}\ 3+\sqrt{5}\\ \textbf{(E)}\ \text{either }3-\sqrt{5}\text{ or }3+\sqrt{5}$


 
 \textbf{Answer: }E



\textbf{Solution (takes advantage of answer choices)}

In order to complete the rectangle area of the oil tank, the chord on the circle must have a length of $4$. Since it's not the diameter $6$, there are two possible outcomes. The only choice that reflects this is $\boxed{\textbf{(E)}}$.

 


	\item (\textit{AHSME 1955}) A three-digit number has, from left to right, the digits $h, t$, and $u$, with $h>u$. 
When the number with the digits reversed is subtracted from the original number, the units' digit in the difference is 4. 
The next two digits, from right to left, are: 


 $\textbf{(A)}\ \text{5 and 9}\qquad\textbf{(B)}\ \text{9 and 5}\qquad\textbf{(C)}\ \text{impossible to tell}\qquad\textbf{(D)}\ \text{5 and 4}\qquad\textbf{(E)}\ \text{4 and 5}$


 
 \textbf{Answer: }B



\textbf{Solution}

We can set up the subtraction like this:
 \[\text{ h t u}\]
 \[- \text{u t h}\]
 \[-------\]
 \[\text{? ? 4}\]Since $u < h$, we need to borrow the one from the tens column. Since the result of the tens column is 0, the taken 1 would result in the tens digit being 9:

 
 \[\text{ h t u}\]
 \[- \text{u t h}\]
 \[-------\]
 \[\text{? 9 4}\]We can assign a value for $u$ and $h$, since that doesn't impact the difference, so lets say that $u=3$ and $h=9$.
 \[\text{ 9 t 3}\]
 \[- \text{3 t 9}\]
 \[-------\]
 \[\text{? 9 4}\]Since $t < t+1$, we can subtract one from the hundred's digit:
 \[\text{ 9 t 3}\]
 \[- \text{3 t 9}\]
 \[-------\]
 \[\text{5 9 4}\]Out difference is $594$. Therefore, the next two digits, from right to left, are $\boxed{\textbf{(B) } \text{9 and 5}}$.

 


	\item (\textit{AHSME 1955}) Four positive integers are given. Select any three of these integers, find their arithmetic average, 
and add this result to the fourth integer. Thus the numbers $29, 23, 21$, and $17$ are obtained. One of the original integers is: 


 $\textbf{(A)}\ 19 \qquad \textbf{(B)}\ 21 \qquad \textbf{(C)}\ 23 \qquad \textbf{(D)}\ 29 \qquad \textbf{(E)}\ 17$


 
 \textbf{Answer: }B



\textbf{Solution}

Define numbers $a, b, c,$ and $d$ to be the four numbers. In order to satisfy the following conditions, the system of equation should be constructed. (It doesn't matter which variable is which.)
 \[\frac{a+b+c}{3}+d=29\]
 \[\frac{a+b+d}{3}+c=23\]
 \[\frac{a+c+d}{3}+b=21\]
 \[\frac{b+c+d}{3}+a=17\]Adding all of the equations together, we get: $2(a+b+c+d)=90$. This means that $a+b+c+d=45$.

 We can determine that $\frac{a+b+c}{3}+d+16=a+b+c+d$. This, with some algebra, means that $\frac{1}{3}(a+b+c)=8$. $d$ must be $\boxed{\textbf{(B)} 21}$.

 


	\item (\textit{AHSME 1955}) If $y=x^2+px+q$, then if the least possible value of $y$ is zero $q$ is equal to: 


 $\textbf{(A)}\ 0\qquad\textbf{(B)}\ \frac{p^2}{4}\qquad\textbf{(C)}\ \frac{p}{2}\qquad\textbf{(D)}\ -\frac{p}{2}\qquad\textbf{(E)}\ \frac{p^2}{4}-q$


 
 \textbf{Answer: }B



\textbf{Solution}

The least possible value of $y$ is given at the $y$ coordinate of the vertex. The $x$- coordinate is given by 
 \[\frac{-p}{(2)(1)} = \frac{-p}{2}\] Plugging this into the quadratic, we get
 \[y = \frac{p^2}{4} - \frac{p^2}{2} + q\]
 \[0 = \frac{p^2}{4} - \frac{2p^2}{4} + q\]
 \[0 = \frac{-p^2}{4} + q\]
 \[q = \frac{p^2}{4} = \boxed{B}\]

 


	\item (\textit{AHSME 1955}) The fractions $\frac{ax+b}{cx+d}$ and $\frac{b}{d}$ are unequal if: 


 $\textbf{(A)}\ a=c=1, x\neq 0\qquad\textbf{(B)}\ a=b=0\qquad\textbf{(C)}\ a=c=0\\ \textbf{(D)}\ x=0\qquad\textbf{(E)}\ ad=bc$


 
 \textbf{Answer: }A



\textbf{Solution}

We can implement each of these options on the existing two fractions:$\textbf{(A)}$: $\frac{x+b}{x+d}$ vs $\frac{b}{d}$. Since $x \neq 0$, we can see that any remaining value of $x$ will make the two unequal. The answer is $\boxed{\textbf{(A)}}$

 $\textbf{(B)}$: $\frac{0}{d}$ vs $\frac{0}{d}$ Both fractions will equal zero $(d \neq 0)$, so they are equivalent.

 $\textbf{(C)}$: $\frac{b}{d}$ vs $\frac{b}{d}$ The two fractions are equal. 

 $\textbf{(D)}$: $\frac{b}{d}$ vs $\frac{b}{d}$ The two fractions are equal.

 $\textbf{(E)}$: This value will vary, but it doesn't take long to figure out a counterexample.

 


	\item (\textit{AHSME 1955}) A train traveling from Aytown to Beetown meets with an accident after $1$ hr. It is stopped for $\frac{1}{2}$ hr., 
after which it proceeds at four-fifths of its usual rate, arriving at Beetown $2$ hr. late. 
If the train had covered $80$ miles more before the accident, it would have been just $1$ hr. late. 
The usual rate of the train is: 


 $\textbf{(A)}\ \text{20 mph}\qquad\textbf{(B)}\ \text{30 mph}\qquad\textbf{(C)}\ \text{40 mph}\qquad\textbf{(D)}\ \text{50 mph}\qquad\textbf{(E)}\ \text{60 mph}$


 
 \textbf{Answer: }A



\textbf{Solution}

Let $v$ be the usual speed of the train in miles per hour, $d$ be the total distance from Aytown to Beetown, and $T = \frac{d}{v}$ be the normal travel time in hours.

 \textbf{Scenario 1:}The train travels at speed $v$ for $1$ hour, stops for $\frac{1}{2}$ hour, and covers the remaining distance $d - v$ at speed $\frac{4}{5}v$. Since the train arrives $2$ hours late, the total time is $T + 2$:
 \[1 + \frac{1}{2} + \frac{d - v}{\frac{4}{5}v} = T + 2\]

 Since $\frac{d - v}{v} = \frac{d}{v} - 1 = T - 1$, we substitute this into the equation:
 \[\frac{3}{2} + \frac{5}{4}(T - 1) = T + 2\]
 \[\frac{3}{2} + \frac{5}{4}T - \frac{5}{4} = T + 2\]
 \[\frac{1}{4} + \frac{5}{4}T = T + 2\]
 \[\frac{1}{4}T = \frac{7}{4} \implies T = 7 \text{ hours}\]

 \textbf{Scenario 2:}If the accident had occurred $80$ miles further, the train would have traveled $1 + \frac{80}{v}$ hours at normal speed before the $\frac{1}{2}$ hour delay. The remaining distance $d - v - 80$ is then covered at speed $\frac{4}{5}v$. Since the train is now $1$ hour late, the total time is $T + 1$:
 \[\left(1 + \frac{80}{v}\right) + \frac{1}{2} + \frac{d - v - 80}{\frac{4}{5}v} = T + 1\]

 We can rewrite the remaining distance term using $\frac{d - v}{v} = T - 1$:
 \[\frac{3}{2} + \frac{80}{v} + \frac{5}{4}\left(T - 1 - \frac{80}{v}\right) = T + 1\]

 Expanding the terms gives:
 \[\frac{3}{2} + \frac{80}{v} + \frac{5}{4}(T - 1) - \frac{100}{v} = T + 1\]

 Substitute $T = 7$:
 \[\frac{3}{2} + \frac{80}{v} + \frac{5}{4}(6) - \frac{100}{v} = 8\]
 \[\frac{3}{2} + \frac{15}{2} - \frac{20}{v} = 8\]
 \[9 - \frac{20}{v} = 8\]
 \[\frac{20}{v} = 1 \implies v = 20\]

 Thus, the usual rate of the train is $\mathbf{20 \text{ mph}}$, which corresponds to $\text{(A)}$.

 


 


	\item (\textit{AHSME 1955}) If $a, b$, and $c$ are positive integers, the radicals $\sqrt{a+\frac{b}{c}}$ and $a\sqrt{\frac{b}{c}}$ are equal when and only when: 


 $\textbf{(A)}\ a=b=c=1\qquad\textbf{(B)}\ a=b\text{ and }c=a=1\qquad\textbf{(C)}\ c=\frac{b(a^2-1)}{a}\\ \textbf{(D)}\ a=b\text{ and }c\text{ is any value}\qquad\textbf{(E)}\ a=b\text{ and }c=a-1$


 
 \textbf{Answer: }C



\textbf{Solution}

Here, we are told that the two quantities are equal. 

 Squaring both sides, we get: $a+\frac{b}{c}=a^2*\frac{b}{c}$.

 Multiply both sides by $c$: $ac+b=ba^2$.

 Looking at all answer choices, we can see that $ac=ba^2-b=b(a^2-1)$. 

 This means that $c=\frac{b(a^2-1)}{a}$, and this is option C. Therefore chose $\boxed{C}$ as your answer.

 


	\item (\textit{AHSME 1955}) The pairs of values of $x$ and $y$ that are the common solutions of the equations $y=(x+1)^2$ and $xy+y=1$ are: 


 $\textbf{(A)}\ \text{3 real pairs}\qquad\textbf{(B)}\ \text{4 real pairs}\qquad\textbf{(C)}\ \text{4 imaginary pairs}\\ \textbf{(D)}\ \text{2 real and 2 imaginary pairs}\qquad\textbf{(E)}\ \text{1 real and 2 imaginary pairs}$


 
 \textbf{Answer: }E



\textbf{Solution}

Graph the two equations. $y=(x+1)^2$ produces a parabola with vertex $(-1, 0),$ while $xy+y=1$ produces a hyperbola. They intersect at the point $(0, 1).$

 This makes one real solution. Only one answer choice has 1 real solution listed, so the answer is $\boxed{\textbf{(E)}}.$

 


	\item (\textit{AHSME 1955}) In circle $O$ chord $AB$ is produced so that $BC$ equals a radius of the circle. $CO$ is drawn and extended to $D$. 
$AO$ is drawn. Which of the following expresses the relationship between $x$ and $y$?


 
                    \begin{figure}[H]
                        centering
                        \includegraphics[width=0.5\textwidth]{Problem/AHSME_I_1955_44_Problem_Figure_1.png}
                    \end{figure}
                    


 

$\textbf{(A)}\ x=3y\\ \textbf{(B)}\ x=2y\\ \textbf{(C)}\ x=60^\circ\\ \textbf{(D)}\ \text{there is no special relationship between }x\text{ and }y\\ \textbf{(E)}\ x=2y\text{ or }x=3y\text{, depending upon the length of }AB$


 
 \textbf{Answer: }A



\textbf{Solution}

Let $r$ be the radius of circle $O$. 

 Since $OB$ is a radius, $OB = r$. We are given that $BC$ equals the radius, so $BC = r$ as well. Therefore, $\triangle OBC$ is isosceles with $OB = BC$. 

 In $\triangle OBC$, the angles opposite the equal sides are equal:
 \[\angle BOC = \angle BCO = y\]

 The exterior angle of $\triangle OBC$ at vertex $B$ is:
 \[\angle ABO = \angle BOC + \angle BCO = y + y = 2y\]

 Since $OA$ and $OB$ are both radii of circle $O$, $OA = OB = r$, making $\triangle OAB$ isosceles with $OA = OB$. Therefore:
 \[\angle OAB = \angle ABO = 2y\]

 Now consider $\triangle AOC$. The exterior angle at vertex $O$ is $x = \angle AOD$. By the exterior angle theorem on $\triangle AOC$:
 \[x = \angle OAC + \angle ACO\]

 Since $\angle OAC = \angle OAB = 2y$ and $\angle ACO = \angle BCO = y$, we have:
 \[x = 2y + y = 3y\]

 Thus, $x = 3y$, which corresponds to $\mathbf{(A)}$.

 


	\item (\textit{AHSME 1955}) Given a geometric sequence with the first term $\neq 0$ and $r \neq 0$ and an arithmetic sequence with the first term $=0$. 
A third sequence $1,1,2\ldots$ is formed by adding corresponding terms of the two given sequences. 
The sum of the first ten terms of the third sequence is: 


 $\textbf{(A)}\ 978\qquad\textbf{(B)}\ 557\qquad\textbf{(C)}\ 467\qquad\textbf{(D)}\ 1068\\ \textbf{(E)}\ \text{not possible to determine from the information given}$


 
 \textbf{Answer: }A



\textbf{Solution}

Let our geometric sequence be $a,ar,ar^2$ and let our arithmetic sequence be $0,d,2d$. We know that 
 \[\begin{cases} a+0=1 \\ ar+d=1\\ ar^2+2d=2\end{cases}\]This implies that $a=1$, hence $r+d=1$ and $r^2+2d=2$. We can rewrite $r+d=1$ as $d = 1-r$, plugging this into $r^2+2d=2$, we get $r^2+2-2r = 2$, we can simplify this to get $r(r-2)=0$, so $r=0$ or $2$. But since $r\neq 0$, $r=2$ and $d=-1$. So our two sequences are $\{1-n\}$ and $\{2^{n-1}\}$, which means the third sequence will be 
 \[\{2^{n-1}+n+1\}=\{1,1,2,5,12,27,\dots\}\]Calculating the sum of the first 10 terms and adding them up yields 978, hence our answer is $\fbox{A}$.

 


	\item (\textit{AHSME 1955}) The graphs of $2x+3y-6=0, 4x-3y-6=0, x=2$, and $y=\frac{2}{3}$ intersect in: 


 $\textbf{(A)}\ \text{6 points}\qquad\textbf{(B)}\ \text{1 point}\qquad\textbf{(C)}\ \text{2 points}\qquad\textbf{(D)}\ \text{no points}\\ \textbf{(E)}\ \text{an unlimited number of points}$


 
 \textbf{Answer: }B



\textbf{Solution}

We first convert each of the lines into slope-intercept form ($y = mx + b$):

 $2x+3y-6=0 ==> 3y = -2x + 6 ==> y = -\frac{2}{3}x + 2$

 $4x - 3y - 6 = 0 ==> 4x - 6 = 3y ==> y = \frac{4}{3}x - 2$

 $x = 2$ stays as is.

 $y = \frac{2}{3}$ stays as is

 We can graph the four lines here:\href{https://www.desmos.com/calculator/ukceuatbb2}{[1]}

 When we do that, the answer turns out to be $\boxed{\textbf{(B)} \text{1 point}}$.

 


	\item (\textit{AHSME 1955}) The expressions $a+bc$ and $(a+b)(a+c)$ are: 


 $\textbf{(A)}\ \text{always equal}\qquad\textbf{(B)}\ \text{never equal}\qquad\textbf{(C)}\ \text{equal whenever }a+b+c=1\\ \textbf{(D)}\ \text{equal when }a+b+c=0\qquad\textbf{(E)}\ \text{equal only when }a=b=c=0$


 
 \textbf{Answer: }C



\textbf{Solution}

Using the FOIL method, we see that $(a+b)(a+c) = a^2 + ab + ac + bc.$ We want to solve 
 \[a + bc = a^2 + ab + ac + bc\]
 \[a = a^2 + ab + ac\]
 \[a((a + b + c) - 1) = 0\]These expressions are $\boxed{\text{(C) equal whenever }  a + b + c = 1.}$

 


	\item (\textit{AHSME 1955}) Given $\triangle$ $ABC$ with medians $AE, BF, CD$; $FH$ parallel and equal to $AE$; $BH and HE$ are drawn; $FE$ extended meets $BH$ in $G$. 
Which one of the following statements is not necessarily correct? 


 $\textbf{(A)}\ AEHF\text{ is a parallelogram}\qquad\textbf{(B)}\ HE=HG\\ \textbf{(C)}\ BH=DC\qquad\textbf{(D)}\ FG=\frac{3}{4}AB\qquad\textbf{(E)}\ FG\text{ is a median of triangle }BFH$


 
 \textbf{Answer: }B



\textbf{Solution}

We evaluate each statement using properties of triangles and parallelograms:

 
                \begin{itemize}
                    \item \textbf{Statement (A):} Quadrilateral $AEHF$ has opposite sides $FH$ and $AE$ that are given to be parallel and equal in length ($FH \parallel AE$ and $FH = AE$). By definition, $AEHF$ is a parallelogram. Thus, \textbf{(A)} is correct.
                \end{itemize}


                
                \begin{itemize}
                    \item \textbf{Statement (C) and (D):} Since $AEHF$ is a parallelogram, $AF \parallel EH$ and $AF = EH$. Since $F$ is the midpoint of $AC$, $AF = FC = EH$, which means $FEHC$ is also a parallelogram. Thus, $FE = CH$.
                \end{itemize}


                Since $F$ and $E$ are midpoints of $AC$ and $BC$, $FE$ is a midline of $\triangle ABC$, so $FE \parallel AB$ and $FE = \frac{1}{2}AB$. 

 In $\triangle BHC$, line $FE$ extended passes through $E$ (the midpoint of $BC$) and is parallel to $CH$. By the converse of the Midpoint Theorem, this line intersects $BH$ at its midpoint $G$, so $EG = \frac{1}{2}CH = \frac{1}{2}FE = \frac{1}{4}AB$. 

 Therefore, $FG = FE + EG = \frac{1}{2}AB + \frac{1}{4}AB = \frac{3}{4}AB$, so \textbf{(D)} is correct.

 
                \begin{itemize}
                    \item \textbf{Statement (E):} Since $G$ is the midpoint of $BH$, segment $FG$ connects vertex $F$ to the midpoint of the opposite side in $\triangle BFH$. Thus, $FG$ is a median of $\triangle BFH$, making \textbf{(E)} correct.
                \end{itemize}


                
                \begin{itemize}
                    \item \textbf{Statement (B):} $HE = AF = \frac{1}{2}AC$, whereas $HG = \frac{1}{2}BH$. Since $AC$ and $BH$ are not equal in a general triangle, $HE = HG$ is not necessarily correct.
                \end{itemize}


                Thus, the statement that is not necessarily correct is $\textbf{(B)}$.

 


	\item (\textit{AHSME 1955}) The graphs of $y=\frac{x^2-4}{x-2}$ and $y=2x$ intersect in: 


 $\textbf{(A)}\ \text{1 point whose abscissa is 2}\qquad\textbf{(B)}\ \text{1 point whose abscissa is 0}\\ \textbf{(C)}\ \text{no points}\qquad\textbf{(D)}\ \text{two distinct points}\qquad\textbf{(E)}\ \text{two identical points}$


 
 \textbf{Answer: }C



\textbf{Solution}

We can simplify $y=\frac{x^2-4}{x-2}$ in order to help solve this problem:  $y=\frac{x^2-4}{x-2}=\frac{(x+2)(x-2)}{x-2}=x+2$.

 In order to find solutions, we have to have $x+2=2x$. $x=2$, but this turns out to be an extraneous solution, as, when put into the two equations, one of them turns out an undefined value. The answer is $\boxed{(C)}$.

 


	\item (\textit{AHSME 1955}) In order to pass $B$ going $40$ mph on a two-lane highway, $A$, going $50$ mph, must gain $30$ feet. 
Meantime, $C, 210$ feet from $A$, is headed toward him at $50$ mph. If $B$ and $C$ maintain their speeds, 
then, in order to pass safely, $A$ must increase his speed by: 


 $\textbf{(A)}\ \text{30 mph}\qquad\textbf{(B)}\ \text{10 mph}\qquad\textbf{(C)}\ \text{5 mph}\qquad\textbf{(D)}\ \text{15 mph}\qquad\textbf{(E)}\ \text{3 mph}$


 
 \textbf{Answer: }C



\textbf{Solution}

Let $V_A, V_B, V_C$ be $A, B, C$'s velocity, respectively. We want to pass $B$ before we collide with $C$. Since $A$ and $B$ are going in the same direction and $V_A>V_B$, $A$ will pass $B$ in $\frac{30\mathrm{ft}}{V_A-V_B}$ time. Since $A$ and $C$ are going in opposite directions, their relative velocity is $V_A+V_C$, so the amount of time before $A$ will collide with $C$ is given by $\frac{210\mathrm{ft}}{V_A+V_C}$. We want to pass $B$ before we collide with $C$, so $V_A$ must satisfy the inequality $\frac{30\mathrm{ft}}{V_A-V_B}<\frac{210\mathrm{ft}}{V_A+V_C}$. We can eliminate all the units, simplifying the inequality to $\frac{1}{V_A-V_B} < \frac{7}{V_A+V_C}$. Solving this and substituting our known values of $V_B$ and $V_C$ yields $330<6V_A$, so $A$ must increase his speed by $\boxed{\textbf{(C) \ } 5 }$ miles per hour.

 


\end{enumerate}